\documentclass{subfiles}
\begin{document}
    \(\np\)-complete problems are too important to be merely ignored because a polynomial solution 
        may not exist. It is for this reason that in most cases an algorithm that computes 
        a near-optimal solution in polynomial time is good enough.
        Still, we need a way to determine how far such an algorithm is from finding the optimal solution.
        A simple solution is considering the approximation ratio \(\rho(n)\).

    Let \(\a\) be an approximation algorithm for a \(\np\)-complete problem,
        and let \(C\) be the cost of the solution found by \(\a\).
        If \(C^{*}\) is the cost of the optimal solution,
        we say that \(\a\) has an approximation ratio \(\rho(n)\) if 
        \[
            \max \left\lparen \frac{C}{C^{*}}, \frac{C^{*}}{C} \right\rparen \le \rho(n)
        \]
        An approximation algorithm with approximation ratio \(\rho(n)\) is also called 
        \emph{\(\rho(n)\)-approximation algorithm}.

    In some cases we can find approximation algorithms that achive better approximation 
        at the cost of more execution time. This kind of algorithms deserve a name on its own.
        We say that an algorithms is an \emph{approximation scheme} if,
        for any \(\epsilon > 0\), the scheme is a \((1 + \epsilon)\)-approximation algorithm.
        We say that an approximation scheme is a \emph{polynomial-time approximation scheme} if 
        for all \(\epsilon > 0\), the schema runs in polynomial time in the size \(n\) of its input.
        We say that an approximation scheme is a \emph{fully} polynomial-time approximation algorithm 
        if it rans in polynomial in both \(\sfrac{1}{\epsilon}\) and \(n\) of its input.
    
    \subsubsection{TSP: an approximation under the triangular inequality}
    \subfile{../subsub/5.1.1 - TSP with triangular inequality}
\end{document}
